% part: intuitionistic-logic
% chapter: introduction
% section: axiomatic-derivations

\documentclass[../../../include/open-logic-section]{subfiles}

\begin{document}

\olfileid{int}{int}{axd}

\olsection{公理\usetoken{P}{derivation}}

直觉主义命题逻辑的公理推导系统是概念上最简单、历史上最早出现的推导系统。其运作方式与经典命题逻辑相同。

\begin{defn}[可推导性]
如果 $\Gamma$ 是 $\Lang L$ 的一个公式集合，那么从 $\Gamma$ 出发的一个\emph{推导}是一个有限的公式序列 $!A_1$, \dots,~$!A_n$，其中对于每个 $i \le n$，以下条件之一成立：
\begin{enumerate}
\item $!A_i \in \Gamma$；或者
\item $!A_i$ 是一条公理；或者
\item $!A_i$ 由满足 $j < i$ 且 $k < i$ 的某个 $!A_j$ 和 $!A_k$ 通过假言推理得到，即 $!A_k \ident !A_j \lif !A_i$。
\end{enumerate}
\end{defn}

\begin{defn}[公理]
直觉主义命题逻辑的\emph{公理}集合 $\PAx$ 由所有具有以下形式的公式组成：
\begin{align}
  & (!A \land !B) \lif !A \ollabel{ax:land1}\\
  & (!A \land !B) \lif !B \ollabel{ax:land2}\\
  & !A \lif (!B \lif (!A \land !B)) \ollabel{ax:land3}\\
  & !A \lif (!A \lor !B) \ollabel{ax:lor1}\\
  & !A \lif (!B \lor !A) \ollabel{ax:lor2}\\
  & (!A \lif !C) \lif ((!B \lif !C) \lif ((!A \lor !B) \lif !C)) \ollabel{ax:lor3}\\
  & !A \lif (!B \lif !A) \ollabel{ax:lif1}\\
  & (!A \lif (!B \lif !C)) \lif ((!A \lif !B) \lif (!A \lif !C)) \ollabel{ax:lif2}\\
  & \lfalse \lif !A \ollabel{ax:lfalse1}
\end{align}
\end{defn}

\begin{defn}[可推导性]
公式~$!A$ \emph{可从} $\Gamma$ \emph{推导}，记作 $\Gamma \Proves !A$，是指存在一个从 $\Gamma$ 出发且以 $!A$ 结尾的推导。
\end{defn}

\begin{defn}[定理]
公式~$!A$ 是\emph{定理}，是指存在一个从空集出发且以~$!A$ 结尾的推导。如果 $!A$ 是定理，我们写作 $\Proves !A$；否则写作 $\Proves/ !A$。
\end{defn}

\begin{prop}
  如果在直觉主义逻辑的公理推导系统中 $\Gamma \Proves !A$，则在经典逻辑中亦有 $\Gamma \Proves !A$。特别地，如果 $!A$ 是直觉主义定理，那么它也是经典定理。
\end{prop}

\begin{proof}
  每条直觉主义公理也是经典公理，因此直觉主义逻辑中的每个推导也是经典逻辑中的推导。
\end{proof}

\end{document}